\chapter{\eh{link} Capa 2 --- Enlace de Datos (Data Link Layer)}

\section{\eh{thinking-face} ¿Qué hace la capa de enlace?}

\begin{concepto}
La capa de enlace mueve \textbf{tramas (frames)} entre dos dispositivos
\textit{directamente conectados} en la misma red local. Usa \textbf{direcciones
MAC} para identificar quién envía y quién recibe dentro de la LAN.

\emoji{light-bulb} \textbf{Analogía:} Si las IPs son como las direcciones
postales que funcionan en todo el mundo, las MACs son como el número de asiento
en un salón de clases: solo funcionan dentro de ese salón (red local).
Cuando el paquete sale de la red local, las MACs cambian en cada salto de router;
la IP permanece igual.
\end{concepto}

\section{\eh{id-button} Dirección MAC}

\begin{concepto}
La \textbf{dirección MAC} (Media Access Control) es un identificador de
\textbf{48 bits = 6 bytes} grabado en la tarjeta de red (NIC) de fábrica.
Teóricamente única en el mundo, aunque puede modificarse por software.
\end{concepto}

\begin{verbatim}
  MAC: 00:1A:2B:3C:4D:5E
       ├── 3 bytes ──┤├─── 3 bytes ───┤
           OUI              NIC-specific
       (fabricante)      (único del dispositivo)

  Ejemplos de OUI:
    00:00:0C → Cisco Systems
    00:50:56 → VMware
    B8:27:EB → Raspberry Pi Foundation
    3C:22:FB → Apple
\end{verbatim}

Tipos de direcciones MAC:
\begin{itemize}
    \item \textbf{Unicast:} dirección única de un solo dispositivo. Bit 0 del primer byte = 0
    \item \textbf{Multicast:} enviar a un grupo. Bit 0 del primer byte = 1
    \item \textbf{Broadcast:} \texttt{FF:FF:FF:FF:FF:FF} --- todos los dispositivos de la LAN la reciben
\end{itemize}

\section{\eh{incoming-envelope} ARP --- Address Resolution Protocol}

\begin{concepto}
\textbf{ARP} (RFC 826) traduce direcciones IP a direcciones MAC.
Es necesario porque Ethernet solo entiende MACs, pero las aplicaciones
usan IPs. Vive entre las capas 2 y 3.

\emoji{light-bulb} \textbf{Analogía:} Sabes el nombre de tu compañero (IP)
pero necesitas saber en qué asiento está (MAC). ARP es como gritar
al salón: ``\textit{¡Oye, Juan García! ¿Cuál es tu asiento?}''
y esperar que Juan levante la mano y diga el número.
\end{concepto}

Proceso ARP completo:
\begin{enumerate}
    \item PC-A (\texttt{192.168.1.10}) quiere hablar con PC-B (\texttt{192.168.1.20})
          pero no conoce su MAC.
    \item PC-A revisa su \textbf{ARP cache}: si ya lo tiene guardado, lo usa.
          Si no:
    \item PC-A envía \textbf{ARP Request} en \textbf{broadcast}
          (\texttt{FF:FF:FF:FF:FF:FF}):
          ``¿Quién tiene \texttt{192.168.1.20}? Respóndele a \texttt{00:AA:BB:11:22:33}''
    \item Todos los dispositivos de la LAN reciben el ARP Request.
          Solo PC-B (\texttt{192.168.1.20}) responde.
    \item PC-B envía \textbf{ARP Reply} en \textbf{unicast} a PC-A:
          ``Yo tengo \texttt{192.168.1.20}, mi MAC es \texttt{00:CC:DD:44:55:66}''
    \item PC-A guarda la respuesta en su \textbf{ARP table} (cache, TTL $\approx$20 min)
    \item PC-A ya puede enviar tramas directamente a la MAC de PC-B
\end{enumerate}

\begin{cuidado}
\textbf{ARP Spoofing / ARP Poisoning:} ARP no tiene autenticación. Un atacante
puede enviar ARP Replies falsos: ``\textit{La IP del router es mi MAC}'' ---
redirigiendo todo el tráfico por su máquina (ataque \textbf{Man-in-the-Middle}).

Defensa: \textbf{Dynamic ARP Inspection (DAI)} en switches Cisco --- valida
las respuestas ARP contra una tabla de DHCP snooping.
\end{cuidado}

\section{\eh{framed-picture} Trama Ethernet (Ethernet II Frame)}

\begin{verbatim}
Ethernet II Frame:
 ┌───────────┬──────────┬────────┬──────────────────┬─────┐
 │ Dst MAC   │ Src MAC  │  Type  │     Payload      │ FCS │
 │  6 bytes  │ 6 bytes  │ 2 bytes│  46 -- 1500 B    │ 4 B │
 └───────────┴──────────┴────────┴──────────────────┴─────┘

 Type field (EtherType):
   0x0800  →  IPv4
   0x0806  →  ARP
   0x86DD  →  IPv6
   0x8100  →  VLAN tagged (802.1Q)
   0x88CC  →  LLDP (Link Layer Discovery Protocol)
\end{verbatim}

El \textbf{FCS} (Frame Check Sequence) es un \textbf{CRC-32}: el emisor calcula
un hash de la trama y lo agrega al final. El receptor recalcula; si no coincide,
la trama se descarta silenciosamente (la capa superior retransmitirá).

\textbf{MTU} (Maximum Transmission Unit): el payload máximo de Ethernet es
\textbf{1500 bytes}. Paquetes más grandes se \textbf{fragmentan} en la capa 3.
Los ``Jumbo frames'' permiten hasta 9000 bytes en redes configuradas para ello
(data centers, iSCSI).

\section{\eh{satellite} Hub vs Switch vs Bridge}

\begin{table}[h!]
\centering
\renewcommand{\arraystretch}{1.4}
\begin{tabular}{llll}
\toprule
\textbf{Característica} & \textbf{Hub} & \textbf{Bridge} & \textbf{Switch} \\
\midrule
Capa OSI        & 1 (Física)     & 2 (Enlace)     & 2 (Enlace) / 3* \\
Inteligencia    & Ninguna        & MAC table (2p) & MAC table (muchos puertos) \\
Reenvío         & Todos los puertos & Por segmento & Solo al puerto destino \\
Dominio colisión & Uno único     & Por segmento   & Por puerto (1 por puerto) \\
Dominio broadcast & Uno          & Uno            & Uno (por VLAN) \\
Estado actual   & Obsoleto       & Reemplazado    & Estándar \\
\bottomrule
\end{tabular}
\caption{Hub vs Bridge vs Switch}
\end{table}

\begin{ejemplo}
\textbf{¿Cómo aprende un switch (MAC learning)?}
\begin{enumerate}
    \item PC-A (MAC \texttt{AA:AA}) envía trama por puerto 1
    \item Switch registra en su MAC table: ``\texttt{AA:AA} → puerto 1''
    \item Destino (MAC \texttt{BB:BB}) es desconocido → \textbf{flooding}:
          manda la trama a todos los puertos excepto el 1
    \item PC-B (MAC \texttt{BB:BB}) en puerto 3 responde
    \item Switch aprende: ``\texttt{BB:BB} → puerto 3''
    \item Comunicaciones futuras: \textbf{forwarding} directo al puerto correcto
\end{enumerate}
\end{ejemplo}

\section{\eh{card-index-dividers} VLANs --- Virtual LANs (IEEE 802.1Q)}

\begin{concepto}
Una \textbf{VLAN} divide lógicamente un switch físico en múltiples switches
virtuales independientes. El tráfico de VLAN 10 es invisible para VLAN 20
y no puede comunicarse con ella sin pasar por un router (capa 3).

\emoji{light-bulb} \textbf{Analogía:} Un edificio con varios departamentos.
Todos comparten el mismo edificio físico (el switch), pero cada departamento
tiene su propia llave y no puede entrar a los otros sin pasar por la
recepción del edificio (el router).

\textbf{Beneficios:} segmentación de seguridad, reducción de broadcast domains,
organización lógica independiente de la ubicación física.
\end{concepto}

El estándar \textbf{IEEE 802.1Q} agrega un tag de \textbf{4 bytes} a la trama
Ethernet entre el campo Source MAC y el EtherType:

\begin{verbatim}
 Trama 802.1Q:
 ┌──────────┬─────────┬────────┬─────┬──────────────┬────────┬─────────┐
 │ Dst MAC  │ Src MAC │ 0x8100 │ PCP │   VLAN ID    │  Type  │ Payload │
 │  6 bytes │ 6 bytes │ 2 bytes│ 3 b │   12 bits    │ 2 bytes│  ...    │
 └──────────┴─────────┴────────┴─────┴──────────────┴────────┴─────────┘
             └──────────────── 802.1Q Tag (4 bytes) ────────┘

 VLAN ID (VID): valores 1–4094
   0    → reservado (sin VLAN)
   1    → VLAN default (no usar para datos)
   4095 → reservado
\end{verbatim}

Tipos de puertos en un switch:
\begin{itemize}
    \item \textbf{Access port:} conecta a un dispositivo final (PC, impresora).
          Solo pertenece a \textit{una} VLAN. Las tramas entran/salen sin tag.
    \item \textbf{Trunk port:} conecta switches entre sí o a un router.
          Transporta tráfico de \textit{múltiples} VLANs con tags 802.1Q.
\end{itemize}

\section{\eh{arrows-counterclockwise} STP / RSTP --- Spanning Tree}

\begin{concepto}
\textbf{STP} (Spanning Tree Protocol, IEEE 802.1D) previene \textbf{loops}
en redes con múltiples switches. Un loop de capa 2 es catastrófico:
los broadcasts se replican infinitamente consumiendo 100\% del ancho de banda
(broadcast storm) hasta colapsar la red.

\emoji{light-bulb} \textbf{Analogía:} Imagina un rumor en la escuela: Juan se
lo cuenta a María, María a Pedro, Pedro a Juan... y el rumor circula para
siempre. STP ``corta'' conexiones redundantes para que no haya loops,
pero las mantiene como respaldo.
\end{concepto}

Proceso STP simplificado:
\begin{enumerate}
    \item \textbf{Elección del Root Bridge:} el switch con menor Bridge ID
          (prioridad + MAC) se convierte en raíz. Por defecto, prioridad = 32768.
    \item \textbf{Root Ports:} cada switch no-raíz elige su mejor camino al Root Bridge.
    \item \textbf{Designated Ports:} en cada segmento, el puerto con mejor camino al Root.
    \item \textbf{Blocking:} los puertos restantes se bloquean (no reenvían datos,
          pero sí escuchan BPDUs para detectar cambios).
\end{enumerate}

\begin{table}[h!]
\centering
\renewcommand{\arraystretch}{1.3}
\begin{tabular}{lll}
\toprule
\textbf{Versión} & \textbf{Convergencia} & \textbf{Nota} \\
\midrule
STP (802.1D)  & 30--50 segundos & Original, lento \\
RSTP (802.1w) & 1--2 segundos  & Rapid STP, estándar actual \\
MSTP (802.1s) & 1--2 segundos  & Multiple STP, una instancia por VLAN \\
\bottomrule
\end{tabular}
\caption{Versiones de Spanning Tree}
\end{table}

\section{\eh{chains} EtherChannel / LACP}

\begin{concepto}
\textbf{EtherChannel} agrupa múltiples enlaces físicos en uno lógico,
multiplicando el ancho de banda y agregando redundancia.

\textbf{LACP} (Link Aggregation Control Protocol, IEEE 802.3ad) es el protocolo
estándar para negociar EtherChannel dinámicamente.
\end{concepto}

\begin{ejemplo}
4 cables Ethernet de 1 Gbps entre dos switches → se ven como
\textbf{un único enlace de 4 Gbps}. Si un cable falla físicamente,
el canal sigue funcionando con 3 Gbps. STP no bloquea ninguno porque
el EtherChannel es un solo enlace lógico.
\end{ejemplo}

\section{\eh{shield} Port Security}

\begin{concepto}
\textbf{Port Security} restringe qué direcciones MAC pueden conectarse a un
puerto del switch. Protege contra MAC flooding y dispositivos no autorizados.
\end{concepto}

Modos de violación:
\begin{itemize}
    \item \textbf{Protect:} descarta tramas de MACs no autorizadas, sin alerta
    \item \textbf{Restrict:} descarta y envía un SNMP trap / log
    \item \textbf{Shutdown:} pone el puerto en \textit{err-disabled state} (requiere intervención manual)
\end{itemize}

\begin{codigo}
\begin{lstlisting}
! Cisco IOS - configurar port security en Fa0/1
interface FastEthernet0/1
 switchport mode access
 switchport access vlan 10
 switchport port-security maximum 2
 switchport port-security mac-address sticky
 switchport port-security violation shutdown
\end{lstlisting}
\end{codigo}
